\clearpage
\chapter*{Errata da exposição n.\textsuperscript{o} 232}
\addcontentsline{toc}{chapter}{Errata da exposição n.\textsuperscript{o} 232}
\src{303}

\noindent\textbf{7. TDTE V: Esquemas de Picard. Teoremas de existência.}

\begin{center}
[Seminário Bourbaki, vol. 14, 1961/62, n.\textsuperscript{o} 232, 19 p.]
\end{center}

\noindent\textbf{Página 7, observação 3.3.} --- A questão levantada aqui foi
resolvida afirmativamente por MUMFORD (ver os comentários à exposição seguinte).

\noindent\textbf{Página 13, observação 5.2.} --- Como assinalo no início da
exposição seguinte, a conjectura de existência formulada aqui é falsa; contudo,
MUMFORD consegue demonstrar, por seus métodos, um teorema ligeiramente mais
fraco.

\noindent\textbf{Página 14, linha 7.} --- Depois de «completo», acrescentar «com
corpo residual algebricamente fechado». O contraexemplo de MUMFORD que acabamos
de mencionar mostra, aliás, que essa restrição é indispensável.

\noindent\textbf{Página 17, observação 6.6.} --- Como assinalamos no início da
exposição seguinte, a questão levantada aqui acaba de ser resolvida
afirmativamente por MURRE.

\noindent\textbf{Página 18, observação 6.7.} --- A questão levantada aqui tem
resposta afirmativa (cf. o último parágrafo dos comentários à exposição
seguinte).

\noindent\textbf{Página 18, observação 6.8, 4\textsuperscript{a} linha.} --- Em vez de
«regular», leia-se «suave sobre $k$».
